一阶\emph{语言}由\emph{常项符号}、\emph{函数符号}和\emph{谓词符号}构成。函数符号和常项符号都带有指定数量的自变元。例如，在\emph{算术语言}中，我们有一个单独的常项符号~$\Obj 0$，一个一元函数符号~$\prime$，两个二元函数符号~$+$ 和~$\times$，以及一个二元谓词符号~$<$。从\emph{变元}以及常项和函数符号出发，我们构造语言的\emph{项}。从语言的项连同其谓词符号以及\emph{等词符号}~$\eq$，我们构造\emph{原子公式}。进而，从原子公式出发，使用逻辑联结词 $\lnot$、$\lor$、$\land$、$\lif$、$\liff$ 以及量词 $\lforall$ 和 $\lexists$，我们构造语言的\emph{公式}。由于在构造项和公式的过程中，我们总是仔细地加上必要的括号，因此每个公式有且仅有一种解读方式。这使得我们可以基于公式的结构进行归纳定义。

公式中变元的出现有时由对应的量词管辖：如果一个变元出现在某个量词的\emph{辖域}内，则它被视为\emph{约束的}，否则是\emph{自由的}。这些概念都有归纳定义，并且我们也归纳地定义了用一个项\emph{代入}公式中一个变元的操作。不含自由变元出现的公式被称为\emph{语句}。